\documentclass[%
  aps, reprint, prl, amsmath,amssymb
]{revtex4-2}

\usepackage{graphicx}% Include figure files
\graphicspath{{figures/}}
\usepackage{orcidlink}
\usepackage{mathtools}
\usepackage{algpseudocode}

% -------------- arXiv-Style Margin Stamp ----------------
\usepackage[style=iso]{datetime2}
\usepackage{FiraMono}
\usepackage{tikz}

% Use shipout/background to ensure it stays behind text/figures
\AddToHook{shipout/background}{
  \begin{tikzpicture}[remember picture, overlay]
    \node [
      rotate=90,
      color=gray!80,
      font=\fontfamily{FiraMono-TLF}\selectfont\footnotesize
    ] at ([xshift=1.2cm]current page.west) {
        WORKING DRAFT \quad $\bullet$ \quad \DTMnow \quad $\bullet$ \quad doi.org/10.17605/OSF.IO/EV8H6 \quad $\bullet$ \quad Brian S. Mulloy \quad $\bullet$ \quad bmulloy@umich.edu
    };
  \end{tikzpicture}
}
% --------------------------------------------------------

\begin{document}

\title{Squeezing the Quantum of Action Determines Information Capacity}

\begin{abstract}
Quantum microstates are operationally indistinguishable inside the quantum of action. Among these infinite possible microstates, the one realized requires the least action to kinematically partition and thermodynamically erase. Least action granulates the continuum. It resolves the coarse-grained configurations of quantum mechanics at low action into the fine-grained configurations of classical mechanics at high action.
\end{abstract}

\author{Brian S. Mulloy\orcidlink{0000-0002-1803-3172}}
\email{bmulloy@umich.edu}

\date{\today}

\maketitle

\section{Introduction}
\begin{equation}
\label{eq:MainIdeaTextWheeler}
\text{Bit} \approx \text{It}
\end{equation}

\begin{equation}
\label{eq:MainIdeaText}
\text{Thermodynamic Erasure} \approx \text{Kinematic Partition}
\end{equation}

\begin{equation}
\label{eq:Hero}
N E_L \tau \le \frac{A}{h}(a_q + a_p)
\end{equation}

\begin{equation}
\label{eq:2DHero}
(N_q + N_p) E_L \tau \le \frac{A}{h}(\pi \delta q \Delta p + \pi \Delta q \delta p)
\end{equation}

\begin{equation}
\label{eq:FullyExpandedHero}
(N_q + N_p) k_B T \ln(2) \tau \le \frac{\pi \hbar}{2} \left( \frac{\Delta q}{\delta q} + \frac{\Delta p}{\delta p} \right)
\end{equation}

\section{The Quantum of Action}

\begin{figure}[t!]
    \centering
    \includegraphics[width=\linewidth]{figures/quantum_of_action.png}
\caption{\textbf{The quantum of action} (schematic, 1D case). By geometrically translating the statistical Heisenberg limit, $\sigma_q \sigma_p \ge \hbar/2$, the classical semi-axes $\Delta q = \sqrt{2}\sigma_q$ and $\Delta p = \sqrt{2}\sigma_p$ establish the relation $\pi \Delta q \Delta p \ge h/2$. In phase space, this is the classical action area, $A \ge h/2$. Confined by these classical boundaries, the interior resolves into a conjugate pair of symplectic quantum blobs, $a_q = h/2$ and $a_p = h/2$, whose sum is the quantum of action, $h$. Specifically, $a_q$ is squeezed by the reciprocal boundary $\delta q = \hbar / \Delta p$ and area $\pi \delta q \Delta p = h/2$, while $a_p$ by the boundary $\delta p = \hbar / \Delta q$ and area $\pi \Delta q \delta p = h/2$. Together, these centered orthogonal ellipses enclose an interior ellipse with quantum action area, $a \le h/2$, evaluated as $\pi \delta q \delta p = \pi \hbar^2/(\Delta q \Delta p)$. This inverse scaling geometrically enforces the correspondence principle: as the classical action area $A \ge h/2$ increases, the symplectic blobs $a_q=h/2$ and $a_p=h/2$ compress into fine cross hairs, shrinking the enclosed quantum action area $a \le h/2$ (see Fig.~\ref{fig:partition_grid}). This symplectic hierarchy establishes the quantum-classical relationships $a = (\pi \hbar)^2 / A$ and $a \le h/2 \le A$.}
\label{fig:quantum_of_action}
\end{figure}

Given phase-space coordinates $(q,p)$, where $q$ is position and $p$ is momentum, the Heisenberg uncertainty principle is $\sigma_q \sigma_p \ge \hbar/2$ \cite{Heisenberg1927}. Geometrically, symplectic mechanics translates this statistical variance into a $1\sigma$ covariance ellipse in phase space, bounding the minimum uncertainty region of the state as $q^2/(2\sigma_q^2) + p^2/(2\sigma_p^2) \le 1$ \cite{DeGosson2006}. The geometric area $A$ of this ellipse is given by $A = \pi (\sqrt{2}\sigma_q) (\sqrt{2}\sigma_p) = 2\pi \sigma_q \sigma_p$.

De Gosson bridges this statistical uncertainty with phase-space geometry. By saturating the Heisenberg bound ($\sigma_q \sigma_p = \hbar/2$), the minimal permissible area emerges as the quantum blob $a_\text{blob}$ \cite{DeGosson2003}:

\begin{equation}
\label{eq:SimpleLanguageOneBlob}
\frac{1}{2} \cdot \text{Quantum of Action} = \text{Quantum Blob}
\end{equation}

\begin{equation}
\label{eq:SymplecticCapacity}
\frac{h}{2} = a_{\text{blob}}
\end{equation}

Protected by Gromov's non-squeezing theorem \cite{Gromov1985}, this $h/2$ area is topologically conserved. Compressing the blob along one axis forces a proportional expansion along its conjugate.

Zurek \cite{Zurek2001} demonstrated that a quantum state develops localized structures as the classical action boundaries, $\Delta p$ and $\Delta q$, compress the conjugate squeezed boundaries, $\delta q = \hbar / \Delta p$ and $\delta p = \hbar / \Delta q$. Expressing these squeezed boundaries as phase-space ellipses yields a conjugate pair of de Gosson blobs (see Fig.~\ref{fig:quantum_of_action}):

\begin{equation}
\label{eq:SimpleLanguageBlobs}
\text{Quantum of Action} = \text{Pair of Conjugate Blobs},
\end{equation}

and their sum is the quantum of action $h=a_q+a_p$ with areas $a_q = \pi \delta q \Delta p = h/2$ and $a_p = \pi \Delta q \delta p = h/2$:

\begin{equation}
\label{eq:Blobs}
h = a_q + a_p
\end{equation}

Furthermore, Zurek showed that quantum interference of the compass state creates structures bounded by these squeezed boundaries, with rectangular area $a_z = \delta q \delta p$ \cite{Zurek2001}. Generalizing to symplectic geometry, the unique maximum-volume shape inscribed within this convex body is the Fritz John ellipsoid \cite{John1948}, shown as the interior quantum ellipse $a \le h/2$ in Fig.~\ref{fig:quantum_of_action}. While the quantum blobs cannot be compressed smaller than de Gosson's $h/2$ limit \cite{DeGosson2022}, squeezing the blobs increases their resolving power forcing this inner geometry to be strictly smaller:

\begin{equation}
\label{eq:JohnEllipsoid}
\pi \delta q \delta p \le \frac{h}{2}
\end{equation}

This establishes classical and action relationships:

\begin{equation}
\label{eq:GeometricUncertaintyPrinciplePlainLanguage}
\text{Quantum Action Area} \propto \frac{1}{\text{Classical Action Area}}
\end{equation}

\begin{equation}
\label{eq:GeometricUncertaintyPrinciple}
a = \frac{(\pi \hbar)^2}{A}
\end{equation}

\begin{equation}
\label{eq:GeometricHierarchy}
a \le \frac{h}{2} \le A
\end{equation}

Substituting the squeezed boundaries, the area of this inscribed ellipse is $a = \pi \delta q \delta p = \pi(\hbar/\Delta p)(\hbar/\Delta q)$, yielding:

\begin{equation}
\label{eq:SqueezedUncertaintyPrinciple}
\delta q \delta p = \frac{\hbar^2}{\Delta q \Delta p}.
\end{equation}

This inverse area scaling geometrically recovers the continuous-variable squeezed states established in quantum optics \cite{Glauber1963, Sudarshan1963, Caves1981}. While the traditional uncertainty principle $\sigma_q \sigma_p \ge \hbar/2$ demands the macroscopic quantum states underlying Schr{\"{o}}dinger's cat paradox \cite{Schrodinger1980Trimmer}, this squeezed uncertainty principle precludes those states and the cat paradox: as classical boundaries expand, they compress the conjugate blobs into fine crosshairs, enforcing the correspondence principle (see the bottom panel of Fig.~\ref{fig:partition_grid}).

\section{Landauer's Principle}

\begin{figure}[htbp!]
    \centering
    \includegraphics[width=\linewidth]{figures/partition_grid.png}
\caption{\textbf{Position densities for the harmonic oscillator}.}
\label{fig:partition_grid}
\end{figure}

Landauer's principle establishes the minimum thermodynamic cost of the thermodynamic erasure of one bit as $E_L = k_B T \ln{2}$ \cite{Landauer1961}.

To express the Landauer erasure of $N$ bits as action, 

\begin{equation}
\label{eq:ErasureActionPlain}
\text{Landauer Action} = \text{Bits} \times \text{Energy} \times \text{Time}
\end{equation}

it is dimensionally necessary to introduce a characteristic erasure time, $\tau$:

\begin{equation}
\label{eq:ErasureActionDefinition}
A_L = N E_L \tau
\end{equation}

We validate this definition at the limit of the minimal quantum phase space, where the pair of quantum blobs degenerate into an overlapping circle with area $h/2$. Saturated by one position bit and one momentum bit, $N = N_q + N_p = 2$, the minimal erasure action becomes:

\begin{equation}
\label{eq:MinimalErasureAction}
2 E_L \tau = \frac{h}{2}
\end{equation}

Solving for $\tau$ recovers the Margolus-Levitin limit \cite{MargolusLevitin1998} for the minimum time to evolve to an orthogonal state ($\tau_\perp$):

\begin{equation}
\label{eq:MLErasureTime}
\tau_{N=2} = \frac{h}{4 E_L} = \tau_\perp
\end{equation}

\begin{equation}
\label{eq:SpeedLimitHero}
N E_L \tau \le \frac{h}{4} \left( \frac{\Delta q}{\delta q} + \frac{\Delta p}{\delta p} \right)
\end{equation}

This demonstrates that the Margolus-Levitin minimal orthogonal time corresponds to the maximal erasure time for a state in phase space specified by a pair of bits $N=N_q+N_p=2$.

Lloyd demonstrated that the number of elementary logical operations $N$ a physical system can perform in time $t$ is bounded by its available energy: $N \le \frac{2 E t}{\pi \hbar}$ \cite{Lloyd2000}. By identifying the erasure of coordinate information as the logical operation of the system, we can map the bits of coordinate precision directly to this computational limit.

Applying this limit to the geometry of phase space, we posit that the total action required to erase $N$ bits of information cannot exceed the classical action area $A$ available to the system:

\begin{equation}
\label{eq:ErasureAndClassicalActionPlain}
\text{Thermodynamic Erasure} \le \text{Kinematic Partitioning}
\end{equation}

\begin{equation}
\label{eq:ErasureAndClassicalAction}
N E_L \tau \le A
\end{equation}

\begin{figure}[t!]
    \centering
    \includegraphics[width=\linewidth]{figures/action_grid.png}
\caption{\textbf{Action dynamics}.}
\label{fig:action_grid}
\end{figure}

Planck demonstrated \cite{Planck1906} that a classical phase space of area $A$ has a configuration capacity of $A/h$. Expanding the quantum of action $h$ into its pair of quantum blobs (Eq.~\ref{eq:Blobs}), we scale the erasure action (Eq.~\ref{eq:ErasureAndClassicalAction}):

\begin{equation}
\label{eq:GeometricActionLimit}
(N_q + N_p) E_L \tau \le \frac{A}{h} (\pi \delta q \Delta p + \pi \Delta q \delta p)
\end{equation}

Where the erasure paths from the least-action microstate $(q_\Upsilon, p_\Upsilon)$ to the blob centers $(q_b, p_b)$ are bounded by $|\epsilon_q| \le \delta q$ and $|\epsilon_p| \le \delta p$.

\subsection{Anatomy of the Erasure Action}

\begin{figure}[t!]
    \centering
    \includegraphics[width=\linewidth]{figures/partition_path.png}
\caption{\textbf{Kinematic partitioning and thermodynamic erasure} (schematic).}
\label{fig:partition_path}
\end{figure}

To understand the physical machinery of Equation~\ref{eq:GeometricActionLimit}, we can isolate the provisioning and expenditure of coordinate bits:

\subsubsection{The Quantum Capacity ($N \propto A/h$)}
The area of the conjugate blob pair is the quantum of action, $h$. Because each quantum blob $a_x = h/2$ can only ever fund the minimal bit budget ($N_q = 1, N_p = 1$), the action required to support a classical bit budget ($N \gg 1$) must be provisioned by the classical scale of the system. The multiplier $A/h$ provides this capacity, scaling the quantum of action $h$ up to the classical action necessary to support the total bits in $N$.

\subsubsection{The Geometric Squeezing ($1/\delta q \propto \Delta p$)}
When a classical boundary expands (e.g., increasing the classical momentum boundary $\Delta p$), symplectic geometry conserves the area of each de Gosson blob at $h/2$. Because this area cannot grow, the blob must physically compress the squeezed position boundary (e.g., $\delta q$) to compensate.

\subsubsection{The Resolving Power ($a \propto 1/A$)}
This symplectic squeezing acts as the resolving power of phase space. To isolate a microstate within shrinking squeezed boundaries ($\delta q \to 0$, $\delta p \to 0$), the phase space must be divided into a strictly larger number of partitions ($N_q \to \infty, N_p \to \infty$). Therefore, classical boundaries enforce this squeezed resolution, which directly dictates the number of bits required for Landauer erasure:

\begin{equation}
\label{eq:ResolvingPowerPlain}
\substack{\text{Classical} \\ \text{Boundary}}
\sim
\substack{\text{Resolving} \\ \text{Power}}
\sim
\substack{\text{Bit} \\ \text{Budget}}
\end{equation}

Isolating the conjugate pairs reveals the causal sequence: the expanding classical boundaries drive the inverse scaling of the squeezed boundaries, which dictates the erasure bit budget:

\begin{equation}
\label{eq:BitBudgetResolutionCoordinate}
\Delta p \sim \frac{1}{\delta q} \sim N_q 
\quad \text{and} \quad
\Delta q \sim \frac{1}{\delta p} \sim N_p 
\end{equation}

When the classical momentum boundary ($\Delta p$) grows, the squeezed position boundary ($\delta q$) must shrink. The universe must therefore spend more bits ($N_q$) to select and erase the state. This changing bit budget dictates the physics we observe: low action affords only enough bits for discrete quantum configurations, while high action supplies the bit budget required to draw smooth classical trajectories.

\section{Least-Action Density}


Standard phase-space formulations of quantum mechanics traditionally employ the Wigner quasi-probability distribution $W(q,p)$ \cite{Wigner1932}. However, because a finite system possesses a finite kinematic bit budget (N), the physically accessible phase space cannot be a continuous manifold; it must be a discretely resolved grid. We instead introduce the least-action density $\Upsilon(q,p)$ as a sum of discrete Iverson brackets \cite{Knuth1992}—mathematical switches that snap to exactly 1 when a coordinate is occupied and 0 otherwise—over the quantum blobs:

\begin{equation}
\label{eq:JointErasure}
\Upsilon(q,p) = \sum_{b} \left[ q = \frac{A}{\pi \hbar} \epsilon_{q,b} \right] \left[ p = \frac{A}{\pi \hbar} \epsilon_{p,b} \right]
\end{equation}

Because the kinematic partitioning evaluates to a discrete coordinate grid, the marginal distributions are obtained by summing over the conjugate variables:

\begin{equation}
\label{eq:Marginals}
N(q) = \sum_{p} \Upsilon(q,p)
\quad \text{and} \quad
N(p) = \sum_{q} \Upsilon(q,p)
\end{equation}

For any phase-space coordinate, $x \in \{q, p\}$, a partitioning function selects the least action microstate $x_\Upsilon$ from among infinite possible microstates $\{x_\mu\}$ inside the squeezed boundary $\delta x$ centered on the blob coordinate $x_b$, yielding a binary sequence $s_x$:

\begin{equation}
\label{eq:PartitionFunction}
(x_\Upsilon, s_x) = \operatorname{PARTITION}(x_b, \delta x)
\end{equation}

Following Bennett's protocol \cite{Bennett1973}, we invert and reverse the binary partition sequence $s_x$ to form the uncomputation path $(\bar{s}_x)^R$, where both paths have bit count $N_x = |s_x| = |(\bar{s}_x)^R|$.

Physical erasure of this sequence uncomputes the microstate back to the ground state ($x_0$). This process expends Landauer action ($\Delta A_L = N_x E_L \tau$) to the environment:

\begin{equation}
\label{eq:ErasureFunction}
\operatorname{ERASE}(s_x, x_\Upsilon, x_b), \quad \text{expends} \quad \Delta A_L = N_x E_L \tau
\end{equation}

The vector $\mathbf{\epsilon} = (\epsilon_q,\epsilon_p)$ is the distance from the least-action state $(q_\Upsilon, p_\Upsilon)$ to the blob center $(q_b, p_b)$. The observed classical coordinates $(q,p)$ are these displacements scaled by the geometric capacity $\frac{A}{\pi \hbar}$:

\begin{equation}
\label{eq:ObservedCoordinates}
(q,p) = \frac{A}{\pi \hbar} (\epsilon_q,\epsilon_p)
\end{equation}

\section{Methods}

\subsection{Dimensional Scaling and Phase-Space Geometry}

To computationally simulate the kinematic partitioning, we map the physical dimensions of phase space to a normalized, dimensionless coordinate system. We establish characteristic reference scales $q_0$ and $p_0$ such that their phase-space area defines the quantum blob in the ground state, $A_0 = \pi q_0 p_0 = h/2$.

Any physical boundary or coordinate can therefore be expressed as a dimensionless ratio, denoted by the $\widetilde{\text{tilde}}$, scaled by its characteristic unit: $\Delta q = \widetilde{\Delta q} q_0$, $\Delta p = \widetilde{\Delta p} p_0$, $\delta q = \widetilde{\delta q} q_0$, and $\delta p = \widetilde{\delta p} p_0$.

The foundation of our topological model rests on four de Gosson ellipsoids. By dividing the physical area of these ellipses by the characteristic capacity $A_0 = h/2$, the physical boundaries reduce to a system of dimensionless geometric constraints:
\begin{enumerate}
    \item \textbf{Classical Action Boundary:} $$\widetilde{\Delta q} \widetilde{\Delta p} \ge 1$$    
    \item \textbf{Squeezed Position Ellipse:} $$\widetilde{\delta q} \widetilde{\Delta p} = 1 \implies \widetilde{\delta q} = 1 / \widetilde{\Delta p}$$
    \item \textbf{Squeezed Momentum Ellipse:} $$\widetilde{\Delta q} \widetilde{\delta p} = 1 \implies \widetilde{\delta p} = 1 / \widetilde{\Delta q}$$
    \item \textbf{Interior Squeezed Ellipse:} $$\widetilde{\delta q} \widetilde{\delta p} = 1 / (\widetilde{\Delta q} \widetilde{\Delta p})$$
\end{enumerate}

\subsection{Semiclassical Harmonic Oscillator}

For the symmetrical harmonic oscillator, the dimensionless classical boundaries expand uniformly: $\widetilde{\Delta q} = \widetilde{\Delta p}$.

The relative action is $A/A_0 = \widetilde{\Delta q} \widetilde{\Delta p}$. For a given action level $A$, the classical boundaries are:
\begin{equation}
\widetilde{\Delta q} = \sqrt{\frac{A}{A_0}} \quad \text{and} \quad \widetilde{\Delta p} = \sqrt{\frac{A}{A_0}}
\end{equation}

Substituting these into the constraints of the quantum blobs $a_q$ and $a_p$ yields:
\begin{equation}
\widetilde{\delta q} = \sqrt{\frac{A_0}{A}} \quad \text{and} \quad \widetilde{\delta p} = \sqrt{\frac{A_0}{A}}
\end{equation}

At the ground state ($A = A_0$), the classical and squeezed boundaries converge: $\widetilde{\Delta q} = \widetilde{\Delta p} = \widetilde{\delta q} = \widetilde{\delta p} = 1$. This topological intersection of four de Gosson ellipsoids marks the ground state, the phase-space boundary where Heisenberg's classical uncertainty area $A$ meets the quantum uncertainty area $a$ and the conjugate pair of quantum blobs $a_q$ and $a_p$, defining a single quantum blob $A=a=a_q=a_p$.

These dimensionless boundaries describe radial phase-space geometry. Because the harmonic oscillator evolves over a closed periodic orbit, projecting this 2D phase space onto a 1D cross-section requires scaling by the time-averaged linear footprint of the orbit. Integrating the absolute displacement over the full angular period yields a geometric projection factor of $2/\pi$. The projected 1D squeezed position boundary is therefore:

\begin{equation}
\widetilde{\delta q}_{\text{orbit}} = \frac{2}{\pi} \sqrt{\frac{A_0}{A}}
\quad \text{and} \quad
\widetilde{\delta p}_{\text{orbit}} = \frac{2}{\pi} \sqrt{\frac{A_0}{A}}
\end{equation}


\begin{equation}
\label{eq:EllipticalToLinearHero}
N E_L \tau \le \frac{\pi}{2} \left( \frac{\Delta q}{\delta q} + \frac{\Delta p}{\delta p} \right)\hbar
\end{equation}

\subsection{Simulation Architecture}

We simulated a total of 990,199,901 individual erasures across 9,901 discrete values of available classical action, sweeping 100,001 independent classical target coordinates ($\tilde{x}_b$) per action level. To ensure uniform data density and avoid quadratic warping at high energies, we performed a linear scan through the classical action space from $A = 1.0 A_0$ to $A = 50.0 A_0$ with a step size of $0.01 A_0$.

To ensure the binary partitioning of the system preserves physical symmetry, the computational engine strictly evaluates the phase space over a normalized reference interval, $\tilde{x} \in [-1.0, 1.0]$. Therefore, the dynamic physical boundaries must be mapped into this relative coordinate space. Dividing the 1D physical squeezed boundary by the expanding classical macroscopic boundary yields the relative algorithmic tolerance required to resolve a single microstate:
\begin{equation}
\widetilde{\delta x} = \frac{\widetilde{\delta x}}{\widetilde{\Delta x}} = \frac{\frac{2}{\pi}\sqrt{A_0/A}}{\sqrt{A/A_0}} = \frac{2}{\pi} \frac{A_0}{A}.
\end{equation}
 
 For each state, the thermodynamic erasure was computed using a custom C++ simulation engine bound to an R-based parallel data pipeline (\texttt{Rcpp} and \texttt{RcppParallel}). To simulate the physical partitioning function $\operatorname{PARTITION}$, the binary phase-space partitioning uses a Signed Stern-Brocot tree \cite{ConcreteMath, Niqui2007}.

Seeded by the blob centers ($\tilde{x}_b$) and the relative algorithmic tolerance ($\widetilde{\delta x}$), the interval is dynamically partitioned into a binary sequence:

\begin{algorithmic}[1]
\Function{PARTITION}{$\widetilde{x_b}, \widetilde{\delta x}$}
    \State $\text{left} \gets (-1, 0)$
    \State $\text{num} \gets 0, \quad \text{den} \gets 1$
    \State $\widetilde{x_\Upsilon} \gets \text{num} / \text{den}$
    \State $\text{right} \gets (1, 0)$
    \State $s_x \gets \text{``''}$
    \While{$|\widetilde{x_\Upsilon} - \widetilde{x_b}| > \widetilde{\delta x}$}
        \If{$\tilde{x}_\Upsilon < \tilde{x}_b$} 
            \State $\text{left} \gets (\text{num}, \text{den})$
            \State $s_x \gets s_x + \text{``1''}$
        \Else 
            \State $\text{right} \gets (\text{num}, \text{den})$
            \State $s_x \gets s_x + \text{``0''}$
        \EndIf
        \State $\text{num} \gets \text{left}_{\text{num}} + \text{right}_{\text{num}}$
        \State $\text{den} \gets \text{left}_{\text{den}} + \text{right}_{\text{den}}$
        \State $\tilde{x}_\Upsilon \gets \text{num} / \text{den}$
    \EndWhile
    \State \Return $\tilde{x}_\Upsilon, s_x$
\EndFunction
\end{algorithmic}

The engine enforces a mathematical condition for compact support; the iterative \texttt{while} loop operates under a closed-set limit, halting the algorithm the moment the relative erasure displacement $|\tilde{\epsilon}_x|$ satisfies the mapped squeezed boundary $\widetilde{\delta x}$. This prevents artificial inflation of the Landauer action and ensures the phase-space volume is bounded by the classical action capacity. 

Once partitioned, the universe must physically erase the state. For example, let $s_x = \mathtt{1101}$. Taking the bitwise complement gives $\bar{s}_x = \mathtt{0010}$, and reversing it yields the uncomputation path $(\bar{s}_x)^R = \mathtt{0100}$. By following $(\bar{s}_x)^R$, the partition sequence $s_x$ is irreversibly destroyed. This erasure dissipates the Landauer action and yields the observable displacement:

\begin{algorithmic}[1]
\Function{ERASE}{$(\bar{s}_x)^R, \widetilde{x_\Upsilon}, \widetilde{x_b}$}
    \State $N_x \gets |(\bar{s}_x)^R|$
    \While{$N_x > 0$}
        \State \Call{ShiftLeft}{$(\bar{s}_x)^R$} 
        \Statex \Comment{Expends action $k_B T \ln{2} \tau$}
        \State $N_x \gets |(\bar{s}_x)^R|$
    \EndWhile
    \State $\widetilde{\epsilon_x} \gets \widetilde{x_\Upsilon} - \widetilde{x_b}$ \Comment{The normalized phase-space displacement}
    \State \Return $\widetilde{\epsilon_x}$
\EndFunction
\end{algorithmic}

The raw output was subsequently processed using multi-threaded routines to calculate the ensemble mean, median, minimum, and maximum sequence lengths, verifying the classical action bound $N E_L \tau \le A$ and extracting the discrete topological nodes from the resulting probability densities.

\begin{figure*}[t!]
    \centering
    \includegraphics[width=\linewidth]{figures/partition_tree.png}
\caption{\textbf{Kinematic partitioning and thermodynamic erasure tree} (schematic).}
\label{fig:partition_tree}
\end{figure*}

\bibliography{main}
\end{document}